\documentclass{article}

\usepackage{amsmath,amssymb}
\usepackage{fullpage}
\usepackage{enumerate}
\usepackage{hyperref}
\usepackage{graphicx}
\graphicspath{{../logos/}}


\begin{document}

\setlength{\tabcolsep}{6pt}
\begin{center} \begin{tabular}{cccc}
	\includegraphics[height=56pt]{SAMF_logo.jpg} &
	\includegraphics[height=56pt]{SAICA_logo.jpg} &
	\includegraphics[height=56pt]{OM_Logo_Stacked_Vignette_on_White_RGB.jpg} &
	\includegraphics[height=56pt]{SAMO2019.png}
\end{tabular} \end{center}


\bigskip


\begin{center}
	\textbf{\Large <LEVEL> Monthly Problem Set Solutions}
\end{center}

\begin{enumerate}

\medskip
\item[1.] % Problem source
Problem statement


\medskip
\item[2.]


\medskip
\item[3.]


\medskip
\item[4.]


\medskip
\item[5.]


\medskip
\item[6.] % Emile Tredoux 2020
Let $E_0,E_1,...,E_{n-2}$ be the set of shared edges in your stack, with $E_i$ being the edge that's shared between the $i^{th}$ and $i+1^{th}$ placed polygons, and with $E_0$ and $E_{n-2}$ being the bottom- and top-most edge respectively.

Note that the edges of a regular $p-gon$ are just rotated by an integer multiple of $\frac{360}{p}$ degrees, therefore edge $E_i$ is just $E_{i-1}$ rotated by some positive integer multiple $a_q < q$ of $\frac{360}{q}$, where $q$ is the number of sides of the polygon shared by $E_i$ and $E_{i-1}$.
Let this multiple be called $q_i$ for $0<i<n-2$

Therefore the angle of inclination of $E_{n-2}$ is just some sum of these $q_i$ mod $180$.
For $E_{n-2}$ to be perfectly flat, this value must equal $0$.

\[ \sum_{i=1}^{n-2} q_i \equiv_{180} 0 \implies {\sum_{i=3}^{n} \frac{a_i}{i} \equiv_{1} 0} \]

Let $p$ be the largest prime less than or equal to $n$.
Note that $\frac{a_p}{p} + \frac{a_i}{i}$ will always have a fractional part with denominator some multiple of $p$ unless $gcd(p,i) \neq 1$, meaning $i \geq 2p$, which by Bertrand's postulate implies that there exists some prime $q$ such that $p < q < 2p$, which is a contradiction.

Therefore no $n$ exist to satisfy the condition in the question statement.

\medskip
\item[7.]


\medskip
\item[8.]


\end{enumerate}

\end{document}
